% Generated draft: requires mathematical and source review.
\documentclass[11pt]{article}
\usepackage[margin=25mm]{geometry}
\usepackage{fontspec}
\setmainfont{Latin Modern Roman}
\usepackage{amsmath,amssymb,mathtools,unicode-math}
\setmathfont{Latin Modern Math}
\usepackage{enumitem,needspace,xurl,hyperref}
\setlength{\parindent}{0pt}
\setlength{\parskip}{5pt}
\setlength{\emergencystretch}{3em}
\newcommand{\N}{\mathbb{N}}
\newcommand{\Z}{\mathbb{Z}}
\newcommand{\Q}{\mathbb{Q}}
\newcommand{\R}{\mathbb{R}}
\newcommand{\C}{\mathbb{C}}
\newcommand{\F}{\mathbb{F}}
\newcommand{\A}{\mathbb{A}}
\newcommand{\PP}{\mathbb{P}}
\newcommand{\E}{\mathbb{E}}
\newcommand{\eps}{\varepsilon}
\newcommand{\dd}{\,\mathrm{d}}
\begin{document}
\section*{UnsolvedMath ID: 20000417}
\textbf{Source problem code:} catalogue entry
\subsection{Definitions and mathematical statement}
J. Getz asks how much middle cohomology of a product \(X=\prod_iX_i\) of modular curves is accounted for by diagonal embeddings of smaller products, and whether additional explicit correspondences are needed. The source also asks about higher-dimensional factors, and D. Ramakrishnan reports Tate classes on four Shimura curves that are Hodge classes but whose cycle representatives were not known in the source. The ground field, levels, coefficient field, precise Tate/Hodge spaces and permitted correspondences are unspecified.

Its two-curve comment credits Ribet with the need for twisting as well as Hecke correspondences. The extracted N\'eron--Severi comparison reverses the quotient: the standard comparison over an algebraic closure, with rational coefficients, is
\[
\frac{\operatorname{NS}(X_1\times X_2)_{\Q}}
{p_1^*\operatorname{NS}(X_1)_{\Q}+p_2^*\operatorname{NS}(X_2)_{\Q}}
\simeq\operatorname{Hom}(J(X_1),J(X_2))_{\Q},
\]
with the usual correspondence/polarization conventions. This correction is explicit rather than treating pullback divisors as a subgroup of the stated Hom group.

The attached aid considers a restricted primitive all-\(H^1\) tensor. For \(2r\) curves over a common number field, choose algebraic projectors \(e_i\) selecting two-dimensional \(\ell\)-adic factors \(V_i\), over a suitable coefficient field, and put
\[
V=(\bigotimes_{i=1}^{2r}V_i)(r),\qquad
T_{\mathrm{geom}}(V)=\bigcup_{K/k\text{ finite}}V^{G_K}.
\]
Assume independent diagonal \(\mathrm{SL}_2\) monodromy on the projective twist-equivalence blocks, with the required finite twists killed after finite extension. These are hypotheses, not conclusions for every product of modular or Shimura curves.

The aid reports that an odd-sized block forces \(T_{\mathrm{geom}}(V)=0\). If every block \(B\) has size \(2r_B\), it reports
\[
\dim T_{\mathrm{geom}}(V)=\prod_B\operatorname{Cat}(r_B),\qquad
\operatorname{Cat}(a)=\frac1{a+1}\binom{2a}{a}.
\]
A basis is indexed by products of noncrossing perfect matchings within the ordered blocks. It reports realizing each paired contraction by a divisor correspondence using Faltings full faithfulness and the idempotent corner \(e_j\operatorname{Hom}(J(X_i),J(X_j))e_i\). Products of the pulled-back divisors then give codimension-\(r\) cycles for the matching basis. If \(T_{\mathrm{pair}}(V)\) is their projected cycle-class span, the reported obstruction quotient is
\[
O(V)=T_{\mathrm{geom}}(V)/T_{\mathrm{pair}}(V)=0.
\]
For one four-factor block, the three pairing products span a two-dimensional plane with one signed Pl\"ucker relation after normalization absorbing K\"unneth Koszul signs. The aid asserts this relation in projected cohomology, not rational equivalence of cycles. This conditional tensor calculation does not settle arbitrary higher-dimensional factors or every historical four-Shimura-curve class.
\subsection{Short English statement}
Determine which middle Tate classes on products of modular or Shimura varieties come from explicit correspondences. The attached aid reports a Catalan matching basis made from divisor products for a projected rank-two tensor under independent \(\mathrm{SL}_2\)-block hypotheses, while retaining the broader source questions.
\subsection{Sources}
\begin{itemize}
\item Local UnsolvedMath registry, permanent ID \(20000417\), exact public code \texttt{catalogue entry}; canonical location \texttt{aim-algebraic-number-theory-notes.json}, notes[68].
\item AIM workshop ``The Tate conjecture,'' Problem 6, attributed to J. Getz; supplied URL \url{https://aimath.org/WWN/tateconjecture/tateconjecture.pdf}. Comments credit Ribet and are attributed to J. Getz, D. Ramakrishnan and D. Ulmer; the four-Shimura-curve comment is attributed to D. Ramakrishnan. The PDF was not consulted here; the literature field is empty.
\item Attached research attempt 1, ``A Catalan correspondence criterion for middle Tate tensors,'' invoking Faltings full faithfulness without full bibliographic citations for it or the block/cycle constructions.
\end{itemize}
\end{document}
